TimesFM~3, released by Google in August 2026, is reported to lead the public forecasting benchmarks. Benchmark data are old and
public, so part of such a lead may reflect exposure during pre-training. We evaluate TimesFM~3 zero-shot under a protocol frozen
before the post-release test data were downloaded, on real data published after the model's release (electricity load and prices, weather, river
flow, exchange and interest rates), on a Wikipedia probe that compares the same pages inside and outside the model's documented
training corpus, on two prospective issues sealed by hash before their outcomes existed, and on 7{,}600 controlled synthetic tasks.
We also reproduce Google's published fev-bench results and read the model's inference code, which has no technical report. The
published results reproduce. On post-release data TimesFM~3 has a 16\% lower MASE and a 29\% lower scaled quantile loss than a
classical ensemble (15\% and 22\% with a re-specified ensemble for river flow), about 5\% lower errors than Chronos-2, which on
hourly data comes from using a longer context, and a lower probabilistic error than TimesFM~2.5; its published benchmark lead over Chronos-2 also
reproduces, and is slightly larger when neither model shares information across series, an option the published Chronos-2 results
use. We find no resolved sign of memorisation; effects of up to about 6\% cannot be excluded. Like every foundation model we ran, TimesFM~3 is worse than the random walk on exchange
rates, interest rates and oil prices (by 11\%). In single synthetic settings (descriptive) it forecasts a spurious decline on
saturating growth and fails after a recent level shift. In synthetic data, multivariate mode helps when a strong cross-series
signal exists and is harmless otherwise; on real data its gain is small. The fixed steps of the inference path contribute little
on average.
